\scanpage{103}

\mspage{\fbox{$1'$}}

Soit $S$ schém., $E$ fibré affine de rg 3 sur $S$, $\Sigma\subset E$ une quadrique affine \inserted{(\unsure{projectiv$^{\mathrm t}$} \deleted{régulière}\inserted{lisse})} -- il revient au $=$ de partir d'un fibré projectif \deleted{\unsure{lisse}} $P$\inserted{$(=\hat E)$} de rang 3, d'une quadrique projective \inserted{lisse} $Q$\inserted{$(=\hat\Sigma)$}$\subset P$, et d'un hyperplan $H$ de $P$ -- en prenant $E=P\setminus H$, $\Sigma=Q\cap E=Q\setminus Q\cap H$. \deleted{On}

[Sauf erreur, il revient encore au $=$ de se donner $Q$ comme schéma (projectif lisse) sur $S$ qui soit une forme de $P^1\times P^1$, ce qui détermine $P$ canoniquement, et de prendre un diviseur \inserted{\deleted{\unsure{$Q_1$}}$X$} de type $(1,1)$ sur $Q$ -- lequel sera défini par un $H$ bien déterminé de $P$\ldots]

On suppose que $\Sigma$ n'est pas un paraboloïde, \inserted{sur aucune fibre} i.e. que $H$ n'est pas tangent à $Q$, ou encore que $\deleted{\unsure{C}}X=H\cap Q=\hat\Sigma\setminus\Sigma$ est lisse i.e. un 1-fibré projectif sur $S$. On peut dire aussi que la donnée de $(E,\Sigma)$ équivaut à $E$, \unsure{sa|son} ext. $\unsure{\mathscr E|\mathcal E}$ de $\mathcal O_S$ par

% archive sheet number: 102
